\documentclass[10pt,a4paper]{amsart}
\usepackage[margin=19mm]{geometry}
\usepackage{amsmath,amssymb,mathtools}
\usepackage[scr=rsfs,cal=euler]{mathalfa}
\usepackage{xfrac}
\usepackage[unicode,hidelinks,bookmarks=false]{hyperref}
\newcommand{\ie}{i.e.\ }
\newcommand{\cf}{cf.\ }
\newcommand{\resp}{respectively\ }
\newcommand{\Subsection}{\subsection}
\providecommand{\nobreakdash}{\nobreak\hbox{-}}
\setlength{\emergencystretch}{2em}
\multlinegap0pt
\usepackage{xcolor}
\usepackage{amssymb}
\newtheorem{definition}{Definition}
\newtheorem{proposition}{Proposition}
\title{\LARGE \bf Optimal Modified Feedback Strategies in LQ Games under Control Imperfections}
\author{Mahdis Rabbani, Navid Mojahed, Shima Nazari}
\date{Source: arXiv:2503.19200v4 (CC BY 4.0)}
\begin{document}
\maketitle

\noindent\emph{Source and licence.} \LARGE \bf Optimal Modified Feedback Strategies in LQ Games under Control Imperfections. Version arXiv:2503.19200v4 (CC BY 4.0), distributed by arXiv under the Creative Commons Attribution 4.0 International licence, http://creativecommons.org/licenses/by/4.0/, as recorded in the arXiv licence metadata for this exact version and shown on the abstract page https://arxiv.org/abs/2503.19200v4. Full licence notice: \texttt{licenses/arxiv-2503.19200-CC-BY-4.0.txt}.\par\smallskip

\noindent\textit{Editorial scope.} Counted unit: the article's proposition prop:first\_order (source0.tex lines 312--331). The article's complete original proof of it is delivered with it (source0.tex lines 333--364, one \texttt{proof} environment of 32 lines). No mathematics is written, repaired or re-derived: the statement and the proof are the article's own lines, in the article's own notation, and no retained line is substituted or re-flowed.  Retained uncounted as the source-local context the proof needs: lines 178--181; lines 262--265; lines 276--279; lines 291--302.  Nothing was omitted from any retained span, so the package carries no omission record.  No retained line contains a citation, so the wrapper carries no bibliography and no bibliography entry.  No retained line contains a figure environment, an image-inclusion command or drawing code, so no image is delivered and the package declares no asset dependency.  Editorial formatting only, in addition: an AMS \texttt{amsart} preamble that restates the article's own declarations and macros used by the retained lines, namely the environments \texttt{definition}, \texttt{proposition} on separate counters, together with the standard packages those lines need.  The retained environments print in this document as Definition 1, Proposition 1 in order of appearance, whereas the article prints the same items with the numbering of its own full document.  The complete native source of the article as distributed in the arXiv e-print for this exact version is bundled unchanged as \texttt{sources/175313/source0.tex}.
\begin{equation}
    \label{eq:LQ cost function}
    J_i = \sum_{k=0}^{N-1} \left( x_k^\top Q_i x_k + u_{i,k}^\top R_i u_{i,k} \right) + x_N^\top Q_{i,N} x_N,
\end{equation}

\begin{equation}
\label{eq:feedback nash definition}
u^\star_{i,k} = -K^\star_{i,k} x_k, \qquad i\in\{1,2\},\ \ k=0,\dots,N-1,
\end{equation}

\begin{equation}
\label{eq:closed-loop Nash dynamic}
x_{k+1}^\star = A_{\mathrm{cl},k} x_k^\star,\qquad x_0^\star \equiv x_0.
\end{equation}

\begin{definition}
\label{def:phi}
For integers $l\ge k\ge j$, define
\begin{equation*}
\Phi(k,j) := \prod_{i=j}^{k-1} A_{\text{cl},i} \;=\; A_{\text{cl},k-1} \cdots A_{\text{cl},j},
\end{equation*}
with the convention $\Phi(j,j):=I$. 
% Then,
% \[
% \Phi(k+1,j)=A_{\text{cl},k}\,\Phi(k,j), \qquad\Phi(k,\ell)\,\Phi(\ell,j)=\Phi(k,j).
% \]
\end{definition}

\begin{proposition}\label{prop:first_order}
Let $\{K^\star_{1,k}, K^\star_{2,k}\}_{k=0}^{N-1}$ be the FNE gains in \eqref{eq:feedback nash definition}. Suppose Player~2 deviates as $u_{2,k}=u^\star_{2,k}+\Delta u_{2,k}$ while Player~1 keeps $u_{1,k}=u^\star_{1,k}=-K^\star_{1,k}x_k$. Then, for $k=1,\dots,N$,
\begin{equation}\label{eq:dx_formula}
\Delta x_k \;=\; \sum_{j=0}^{k-1} \Phi(k,j+1)\,B_2\,\Delta u_{2,j},
\end{equation}
and, with $S_k := Q_1 + K_{1,k}^{\star\top} R_1 K^\star_{1,k}$,
\begin{equation}\label{eq:DJ_linear}
\Delta J_1 \;=\; 2 \sum_{j=0}^{N-1} x_0^\top \Lambda_j \,B_2\,\Delta u_{2,j} \;+\; \mathcal{O}\!\big(\|\Delta u_2\|_2^2\big),
\end{equation}
where
\begin{equation}\label{eq:Lambda_j}
\begin{aligned}
    \Lambda_j \;=\; \sum_{k=j+1}^{N-1} \Phi(k,0)^\top &S_k \,\Phi(k,j+1) \\
    \;&+\; \Phi(N,0)^\top Q_{1,N}\,\Phi(N,j+1),
\end{aligned}
\end{equation}
for $j=0,\dots,N-1$.
Here $\Phi(k,j)$ is as in Definition~\ref{def:phi}, $\Delta u_2:=[\Delta u_{2,0}^\top\,\cdots\,\Delta u_{2,N-1}^\top]^\top$, and
the remainder is with respect to the Euclidean norm $\|\Delta u_2\|_2$.
\end{proposition}

\begin{proof}
Under the FNE
% and defining $x_0^\star \equiv x_0$
, the nominal trajectory is obtained via \eqref{eq:closed-loop Nash dynamic}. With $u_{2,k}=u_{2,k}^\star+\Delta u_{2,k}$ and $u_{1,k}=u_{1,k}^\star$, the actual state satisfies
\begin{equation*}
x_{k+1} \;=\; A_{\text{cl},k} \, x_k + B_2 \, \Delta u_{2,k}.
\end{equation*}
Defining $\Delta x_k := x_k - x_k^\star$ gives
\begin{equation*}\label{eq:dx_recursion}
\Delta x_{k+1} \;=\; A_{\text{cl},k} \, \Delta x_k + B_2 \, \Delta u_{2,k}, \qquad \Delta x_0 = 0,
\end{equation*}
whose iteration yields \eqref{eq:dx_formula} using $\Phi(k+1,j)=A_{\text{cl},k}\,\Phi(k,j)$.

Next, linearizing \eqref{eq:LQ cost function} for $i=1$ around the FNE trajectory $(x_k^\star,u_{1,k}^\star)$, leads to
\begin{equation*}\label{eq:DJ_raw}
\begin{aligned}
    \Delta J_1 \;=\; 2 \sum_{k=0}^{N-1} \big( x_k^{\star\top}& Q_1 \Delta x_k + u_{1,k}^{\star\top} R_1 \Delta u_{1,k} \big)\\
    \;+\; &2 \, x_N^{\star\top} Q_{1,N} \Delta x_N \;+\; \mathcal{O}(\|\Delta x\|_2^2),
\end{aligned}
\end{equation*}
where $\Delta\,J_1 = J_1(x_k,u_{1,k})-J_1(x_k^\star,u_{1,k}^\star)$.

Since $u_{1,k}=-K^\star_{1,k}x_k$, we have $\Delta u_{1,k} = -K^\star_{1,k}\Delta x_k$, hence
\begin{equation}\label{eq:DJ_Sk}
\begin{aligned}
    \Delta J_1 \;=\; 2 \sum_{k=0}^{N-1} x_k^{\star\top} S_k \, \Delta x_k
\;+\; 2 x_N^{\star\top}& Q_{1,N} \Delta x_N \\
\;&+\; \mathcal{O}(\|\Delta x\|_2^2),
\end{aligned}
\end{equation}
where $S_k := Q_1 + K_{1,k}^{\star\top} R_1 K^\star_{1,k}$. Substituting \eqref{eq:dx_formula} into \eqref{eq:DJ_Sk}, defining $\Lambda_j$ as in \eqref{eq:Lambda_j}, and using $x_k^\star=\Phi(k,0)x_0$ gives \eqref{eq:Lambda_j}. As $\Delta x$ is linear in $\Delta u_2$ by \eqref{eq:dx_formula}, the remainder is $\mathcal{O}(\|\Delta u_2\|_2^2)$. This completes the proof.
\end{proof}

\end{document}
